\section*{Appendix}

\subsection*{[1] Estimating the Model Impact Parameters}

\cite{online:orderbook-optimal-liquidation} proposed practical, stylized heuristics to approximate these constants based on readily available market data. This method assumes that trading a specific percentage of the daily volume moves the asset price by exactly one full bid-ask spread.
\begin{itemize}
\item Temporary Impact ($\eta$): Assumes that executing a trade size equal to $1\%$ of the daily volume incurs a temporary impact (slippage) equal to one full bid-ask spread.
$$\eta =\frac{\text{Bid-Ask\ Spread}}{0.01\times \text{Daily\ Volume}}$$

\item Permanent Impact ($\gamma$): Assumes that executing a trade size equal to $10\%$ of the daily volume permanently shifts the equilibrium price by one full bid-ask spread.
$$\gamma =\frac{\text{Bid-Ask\ Spread}}{0.10\times \text{Daily\ Volume}}$$
\end{itemize}

For a more precise institutional setup, calculate $\eta$ and $\gamma$ by running statistical regressions on historical execution data. The \cite{online:orderbook-optimal-liquidation} framework breaks down the realized execution price $S_{t}$ at time step $t$ using the following linear equations
\begin{enumerate}
\item Gather the Data: Collect historical data on the executed orders. For each trading interval, obtain the trade execution volume $(n_{t}$, the trading speed $v_t = n_t / \tau$, the asset's unaffected mid-price $S_{t}^{0}$, and the actual execution price $S_{t}.

\item Estimate Permanent Impact ($\gamma$): Measure the net price change over the entire duration of a metaorder. Run a linear regression without an intercept: 
$$\Delta S_{\text{\text{permanent}}}=\gamma \cdot Q+\epsilon$$
where:
\begin{itemize}
\item $v_t$: trading rate
\item The resulting slope coefficient of regression will be the $\eta$
\end{itemize}
\end{enumerate}

\subsection*{[2] Hamilton Jacobi Bellman Equation}
The continuous time version of the discrete Bellman equation takes the form:
$$-\frac{\partial V}{\partial t} = \min_v \bigg\{ \mathcal{L}(q,v) + \frac{\partial V}{\partial q} f(q,v) \bigg\}$$

The term inside the $\min$ operator is called the Hamiltonian, denoted as $\mathcal{H}(q,v,\frac{\partial V}{\partial q})$

$$\mathcal{H}\left(q,v,\frac{\partial V}{\partial q} \right) = L(q,v) + \frac{\partial V}{\partial q} f(q,v)$$

The HJB equation is
$$-\frac{\partial V}{\partial t} = \min_v\left\{ \mathcal{H} \left(q,v, \frac{\partial V}{\partial q}\right) \right\}$$

To minimize the Hamiltonian, we differentiate $\mathcal{H}$ with respect to $v$ and set the result to $0$. This gives the optimal trading rate $v^*(q,t)$, which specifies the best trading rate for any given state $q$ and time $t$.

The derived $v^*(q,t)$ must respect the physical constraints by clipping the trading rate to the allowed limits $[v_{min}, v_{max}]$.

After finding $v*(q,t)$ we substitute it back into the Hamiltonian to find $H_{min}$ that represents the rate of change of the value function with respect to time.

$$\mathcal{H}_{min}(q,t) = L\big(q, v^*(q,t)\big) + \frac{\partial V}{\partial q} f\big(q, v^*(q,t)\big)$$
The $H_{min}$ represents the change of the value function with respect to time.

The HJB equation is a partial differential equation, and to solve it uniquely, we need a boundary condition. In optimal control problems solved backward in time, this is usually a terminal condition at time $T$.
$$V(q,T) = 0$$

We solve numerically backward from the terminal state. The HJB equation, coupled with the terminal condition, defines the optimal value function. We implement the numerical solution using a backward time iteration and finite difference method:

\begin{itemize}
\item Set $V(q, T)$ using the terminal cost
\item For each time step backward from $T$ to $0$
\begin{itemize}
\item Approximate $\frac{\partial V}{\partial q}$ using finite differences on the current 
$V(q, t)$.
\item Calculate the optimal control 
$v^*(q, t)$ using the derived formula and clip it.
\item Compute the minimum $H_{min}(q, t)$ using $v^*$.
\item Update the value function for the previous time step: 
$V(q, t - \Delta t) \approx V(q, t) - H_{min}(q, t) \Delta t$.
\end{itemize}
\end{itemize}

By iterating this process, we progressively compute 
$V(q, t)$ for earlier times, eventually arriving at 
$V(q, 0)$, which gives the minimum cost from any initial state $q$ at time $t=0$. The sequence of 
$v^*(q, t)$ computed at each step forms the optimal control policy.